\section{\systemname}
\label{sec:method}

\cref{fig:pipeline} shows the three-stage \systemname{} pipeline: collect trajectories, propose trajectory-level patches in parallel, and consolidate them into one portable skill.
We first define the skill-evolution objective, then describe each stage.

\subsection{Skill and Problem Formalization}
\label{sec:method:skill}

A skill is a human-readable directory $\mathcal{S}=(M,\mathcal{R})$, where $M$ is the root \texttt{SKILL.md} and $\mathcal{R}$ contains auxiliary references, scripts, or assets.
The root document stores broadly applicable procedural knowledge, while auxiliary files provide deterministic tools or lower-frequency details.
Let $\pi_\theta$ be a fixed LLM agent using skill $\mathcal{S}$ at inference time, and let $\mathcal{P}(\mathcal{S};\pi_\theta,\mathcal{D})$ denote the pass rate of that agent on task set $\mathcal{D}$.
Given an evolving set $\mathcal{D}_\text{evolve}$ and a disjoint test set $\mathcal{D}_\text{test}$, skill evolution constructs a new skill from evolving-set trajectories without updating $\theta$:
\begin{equation}
    \mathcal{S}^* = \mathcal{E}(\mathcal{S}_0,\mathcal{D}_\text{evolve};\pi_\theta),
    \qquad
    \mathcal{P}(\mathcal{S}^*;\pi_\theta,\mathcal{D}_\text{test})
    > \mathcal{P}(\mathcal{S}_0;\pi_\theta,\mathcal{D}_\text{test}).
\end{equation}
We evaluate two initializations for $\mathcal{S}_0$: a human-expert skill and an LLM-generated draft from parametric knowledge alone.

\subsection{Stage 1: Trajectory Generation}
\label{sec:method:trajectories}

We use a ReAct-style harness \citep{yao2023reactsynergizingreasoningacting}.
For each evolving-set task, the fixed agent runs with $\mathcal{S}_0$ and produces a trajectory $\tau_i$ containing the query, reasoning/tool-use history, final output, and a binary correctness outcome.
The resulting corpus $\mathcal{T}$ is split into failures $\mathcal{T}^-$ and successes $\mathcal{T}^+$.
Trajectory generation is parallel across evolving-set problems, so the trace corpus can be collected independently before patch proposal.
See prompt templates in \cref{app:stage1_prompt}.

\subsection{Stage 2: Parallel Patch Proposal}
\label{sec:method:swarm}

A group of analyst sub-agents independently proposes skill patches from individual trajectories.
Failures are sent to an error analyst $\mathcal{A}^-$, successes to a success analyst $\mathcal{A}^+$, and the analysts read $\mathcal{S}_0$ and propose a patch for it, leading to a patch pool $\mathcal{P}=\mathcal{P}^-\cup\mathcal{P}^+$.

The analyst roles are intentionally asymmetric.
$\mathcal{A}^+$ uses a single-pass workflow to identify reusable behavior patterns from successful trajectories.
$\mathcal{A}^-$ uses a ReAct-style loop that can inspect traces and artifacts, compare outputs against ground truth, and validate candidate fixes before proposing a patch.
Failures that cannot be causally explained are excluded, ensuring $\mathcal{P}^-$ is grounded in verified failure mechanisms rather than log-only guesses \citep{ouyang2026reasoningbankscalingagentselfevolving}.
Both roles are prompted to write concise, actionable patches following skill-writing guidance \citep{anthropic2026skillcreatorconversation}.
Analyst prompt templates and representative example patches are provided in \cref{app:stage2_prompts}.

\subsection{Stage 3: Patch Consolidation}
\label{sec:method:consolidation}

Stage~3 consolidates the patch pool into one coherent update $p^*$ and applies it to $\mathcal{S}_0$.
Patches are merged hierarchically for $L=\lceil \log_{B_\text{merge}}|\mathcal{P}| \rceil$ levels; at each level $\ell$, up to $B_\text{merge}$ patches are synthesized into one patch:
\begin{equation}
    p^{(\ell+1)}
    = \mathcal{M}\!\left(\pi_\theta,\; \mathcal{S}_0,\; \{p_1^{(\ell)}, \ldots, p_{B_\text{merge}}^{(\ell)}\}\right),
    \qquad \ell=0,\ldots,L-1,
\end{equation}
where $\mathcal{M}$ deduplicates, resolves conflicts, and preserves non-overlapping insights.
$\pi_\theta$ itself serves as trajectory generator, analyst, and merge operator, so no external evolution/teacher model is required.
The final $p^*$ is translated into diff-style edits and applied with deterministic guardrails: reject edits to missing files, withhold line-range conflicts, and validate the updated skill format.

The merge also performs inductive generalization.
Because each patch comes from one trajectory, recurring edits across independent patches are evidence of systematic task properties rather than one-off fixes.
$\mathcal{M}$ is therefore instructed to prefer prevalent patterns and discard idiosyncratic edits, producing a compact skill that replaces $\mathcal{S}_0$ and is used directly at inference without retrieval. The merge operator prompt template and an example consolidated patch $p^*$ are given in \cref{app:stage3_prompts}.

% \subsection{Two Evolution Modes}
% \label{sec:method:modes}

\systemname{} supports two modes that cover scenarios with and without expert-written skills: \emph{Skill deepening} starts with a human-written skill, and \emph{Skill creation} starts from a skill generated from LLM parametric knowledge. Both improve the target skill using patches induced from trace analysis.
